\chapter*{前言}




\subsection*{过去十年的因果推断研究与教育}



过去十年间，由于因果推断在生物医学研究、社会科学及科技企业等领域有着广泛应用，相关研究与教育也迎来了蓬勃发展。即便是在十年前，当我还是统计学专业的博士生时，情况也大不相同。那时，因果推断尚未成为统计学领域的主流研究主题，极少有本科或研究生课程开设因果推断相关内容。在统计学学术界，许多人对因果推断的理论基础仍持强烈怀疑态度。许多顶尖统计学家因其在概念层面存在的根本性难题而不愿接受因果推断——这些难题与传统数理统计学的训练方式截然不同。 


因果推断在实证研究中的应用，彻底改变了统计学领域的研究与教育方式。归根结底，统计学并不仅仅关乎抽象理论，更在于解决现实世界中的各类问题。众多才华横溢的研究者纷纷投身于推动因果推断相关知识的发展。许多学生也热切希望掌握最前沿的因果推断理论与方法，从而更有能力去解决各个领域中的实际问题。 



出于学生的需求，我的同事们鼓励我开设一门关于因果推断的课程。起初，我讲授的是一门同时面向政治学与统计学专业研究生的课程；该课程此前由我的前同事Jas Sekhon在加州大学伯克利分校任教多年。后来，我又为本科生和研究生分别开发了相应的课程。在加州大学伯克利分校，“因果推断”课程的编号分别为Stat 156与Stat 256：本科生修读Stat 156，研究生则修读Stat 256。尽管两个班次的学员需要完成不同的作业、阅读任务以及期末项目，但他们使用的讲义完全相同，也共同聆听了由我及我的助教们主讲的课程。  


鉴于我的学生们在技术准备程度上存在较大差异，教学过程中最具挑战性的部分便是如何兼顾本科生与研究生的兴趣需求。一方面，我希望以直观易懂的方式讲解内容，仅要求本科生具备概率论、统计学、线性回归以及逻辑回归方面的基础知识。另一方面，我也希望向研究生们介绍最新的研究课题与成果。本书正是我过去七年努力的成果。  



\subsection*{给授课教师的建议}



本书的正文部分包含29章，附录部分有3章。加州大学伯克利分校采用学期制，每个学期的授课时长均为14周。我无法在一个学期内讲完全部32章的内容。以下是我结合自身教学经验给出的一些建议。


\paragraph*{附录}
我最初安排学生学习正文中的各章节内容，但也要求我的助教们带领学生复习\ref{chapter::basic-prob-append}章和\ref{appendix::basic-linear-regression}章的基础知识。为了促使学生在阅读正文前先复习\ref{chapter::basic-prob-append}至\ref{chapter::SRS-appendix}章的内容，我还在学期初布置了若干源自\ref{chapter::basic-prob-append}至\ref{chapter::SRS-appendix}章的习题。



\paragraph*{第 \ref{part::introduction} 部分 }
第 \ref{chapter::correlationassociation} 章的核心主题是尤尔–辛普森悖论。第 \ref{chapter::potential-outcomes} 章介绍了潜在结果这一概念，它也是整本书的基础。 



\paragraph*{第一部分 \ref{part::rcts}}
对于随机实验相关的内容，不同的研究者和教学者可能会持有截然不同的观点。我曾与许多朋友探讨过目前在\ref{part::rcts}部分中对随机实验的阐述方式究竟有何优劣。由于随机化过程能够消除未测量的混杂因素，因此在随机实验中开展因果推断相对简单。正因如此，部分朋友认为\ref{part::rcts}部分的第\ref{ch::frt-cre}章至第\ref{chapter::bridging}章篇幅过长，至少对于初学因果推断的人来说如此。当我花费整整一个月时间讲解随机实验时，这也让部分学生感到有些失望。

另一方面，我最初是通过\citet{imbens2015causal}所著的书籍接受相关训练的，我一直认为：要想深入理解观察性研究，首先应当掌握随机实验的相关知识。此外，我也是\citet{Neyman:1923}与\citet{Fisher:1935}所开展的经典研究的忠实拥趸。正因如此，\ref{part::rcts}部分的内容充分体现了我个人的学术成长历程以及对统计学领域的个人偏好。当然，其他教学者或许并不希望花费整整一个月时间讲解随机实验，他们也可以选择快速讲解\ref{chapter::stratification-poststratification}、\ref{chapter::mpe}、\ref{chapter::unification-fisher-neyman}以及\ref{chapter::bridging}这几章的内容。


\paragraph*{第 \ref{part::observational-studies}} 部分
第 \ref{part::observational-studies} 部分介绍了不存在未测量混杂因素的观察性研究中的核心概念。观察性研究的四大支柱分别是结果回归、逆倾向得分加权、双重稳健估计以及匹配估计法，这些内容将分别在第 \ref{chapter::observational-studies}、\ref{chapter::pscore-key}、\ref{chapter::doubly-robust} 和 \ref{chapter::matching-obs} 章中详细阐述。
第 \ref{chapter::ATT-and-other} 章与第 \ref{chapter::pscore-unification} 章属于教学中的可选内容。不过第 \ref{chapter::ATT-and-other} 章与第 \ref{chapter::pscore-unification} 章中的研究结果同样颇具价值，因此我有时会选取其中一两个结果进行讲解，也会要求助教在实验课上补充更多内容，同时还会布置这些章节中的习题以鼓励学生自行阅读。 


\paragraph*{第 \ref{part::challenges-os}}部分
第 \ref{part::challenges-os}部分以新颖的方式探讨了观察性研究中的各类基本难题，包括未测量混杂因素以及重叠性问题。不过，鉴于这些问题的复杂性与微妙性，该部分内容远称不上完美。其中，第 \ref{chapter::evalue}、\ref{chapter::sensitivity-ACE}和\ref{chapter::overlap}章为核心内容，而第 \ref{chapter::difficulties-observational-studies}和\ref{sec::rosenbaum-sensitivity-analysis}章则属于选读部分。 


\paragraph*{第 \ref{part::instrumentalvariables}}部分
第 \ref{part::instrumentalvariables}部分探讨了工具变量的概念。其中第 \ref{chapter::iv-experiment}、\ref{chapter::iv-econometric}和 \ref{chapter::iv-frdd}章至关重要，而第 \ref{chapter::iv-inequalities}和 \ref{chapter::iv-mr}章则属于选读内容。


\paragraph*{第\ref{part::post-treatmentvariable}}部分
第\ref{part::post-treatmentvariable}部分涉及一些特殊主题。从某种意义上说，这些内容都是可选的。鉴于Baron与Kenny提出的中介分析方法广受欢迎，讲授第\ref{chapter::mediation}章或许是值得的。 


\paragraph*{ 略过的主题}
本书并未涉及一些常用的计量经济学方法，例如双重差分法、面板数据分析以及合成控制法。授课教师可参考\citet{angrist2008mostly}以了解这些主题的相关内容。
本书假定读者仅具备最基础的概率与统计学知识。由于大多数入门级统计学课程均采用频率学派的观点，即认为未知参数是固定不变的，因此我在本书中也采纳了这一观点，而未涉及贝叶斯学派在因果推断中的应用。实际上，许多因果推断的基本概念均源自贝叶斯学派，这一点从\citet{rubin1978bayesian}便可看出。若读者或学生希望深入了解贝叶斯因果推断的相关内容，建议阅读由\citet{li2023bayesian}撰写的综述性论文。 



\paragraph*{教学助手的帮助}  
我在加州大学伯克利分校授课时，教学助手们给予了极大的帮助。由于我的讲义中无法涵盖所有内容，因此我经常依靠他们来讲解一些技术细节，或是组织实验室中的\ri{R}编程环节。 
 


\paragraph*{部分作业题的解答}
我还准备了大多数理论题的解答。如果您是因果推断课程的授课教师，请与我联系，我将为您提供包含课程详细信息的解答资料。  





\subsection*{给读者与学生的额外建议}



读者和学生们可以先阅读上文中我给授课教师的建议。此外，我还有两条其他建议。 


\paragraph*{作业习题}
本书的每一章都配有作业习题。为了加深理解，尝试解答一些习题十分重要。此外，部分习题还包含了有用的理论结果。即便你没有时间去详细解答这些习题，至少阅读一下题目的描述也会有所助益。 

\paragraph*{推荐阅读}
本书的每一章都列出了推荐阅读内容。如果你想从事因果推断方面的研究，这些推荐的文献可作为相关的背景知识。我在加州大学伯克利分校讲授研究生级别的因果推断课程时，从第一周起直至学期结束，每周都会布置以下论文作为学生的阅读材料：  
\begin{itemize}
\item 
\citet{bickel1975sex}; 
\item 
\citet{holland1986statistics};
\item
\citet{miratrix:2013};
\item
\citet{lin2013};
\item 
\citet{li2018asymptotic};
\item 
\citet{rosenbaum1983central};
\item 
\citet{lunceford2004stratification};
\item
\citet{ding2016sensitivity};
\item
\citet{pearl1995causal};
\item
\citet{angrist1996identification};
\item
\citet{imbens2014instrumental};
\item
\citet{frangakis2002principal}. 
\end{itemize}
许多学生都向我反馈，称阅读上述论文的经历让他们受益匪浅。  
即便您并不打算阅读本书，我也推荐您去读一读那些论文。 




 


\subsection*{本书的特点}


在过去的十年里，已经出版了许多优秀的因果推断相关书籍。其中一些对我产生了深远的影响。  
大学时期，我在网上阅读了\citet{imbens2015causal}的部分初稿章节。它们彻底颠覆了我对统计学原有的认知方式，也帮助我确立了自己对因果推断的研究兴趣。  
\citet{angrist2008mostly}这本书我反复阅读多次，每次重读都能获得新的启发。  
\citet{rosenbaum2002design}、\citet{morgan2015counterfactuals}和\citet{hernan2020causal}则是另外三本由顶尖因果推断研究者撰写的优秀著作。  
在我筹备本书期间，\citet{cunningham2021causal}、\citet{huntington2022effect}、\citet{brumback2022fundamentals}以及\citet{huber2023causal}这四本近期出版的优秀因果推断相关书籍也相继问世。 


得益于我在加州大学伯克利分校的教学经验，本书具备以下特点，相信会深受授课教师、学生及其他读者青睐。 
\begin{itemize}
\item
本书假定读者对因果推断并无太多预备知识，并在附录中回顾了基本的概率与统计学知识。

\item
本书从统计学、生物统计学及计量经济学的角度探讨了因果推断问题，并列举了多个领域的实际应用案例。 

\item
本书借助\ri{R}代码与数据分析来阐释因果推断的相关概念。所有\ri{R}代码及数据集均可在哈佛数据平台上免费获取：\url{https://doi.org/10.7910/DVN/ZX3VEV}

\item
本书包含课后习题，可作为本科生及研究生的教材使用。授课教师亦可向我索取部分习题的解答。 
\end{itemize}




\subsection*{致谢}



我在大学求学期间，北京大学的耿直教授向我介绍了因果推断这一领域。后来在哈佛大学统计系攻读博士学位时，我的论文委员会成员包括卢克·米拉特里克斯教授、蒂尔坦卡尔·达斯古普塔教授以及唐·鲁宾教授。此后，我在哈佛大学陈曾熙公共卫生学院从事流行病学方向的博士后研究，泰勒·范德韦勒教授担任我的指导老师。   


我在伯克利分校统计系的同事们共同营造了一个极具价值且富有成效的研究环境。自我还是一名初级教员起，Bin Yu与Jas Sekhon便一直给予我诸多支持。系里的各位主任，即Deb Nolan、Sandrine Dudoit以及Haiyan Huang，也鼓励我开设“因果推断”这门课程。这段经历对我而言收获颇丰。  


我很幸运能够与众多合作者共事，尤其是Avi Feller、Laura Forastiere、Zhichao Jiang、Fang Han、Fan Li、Xinran Li、Alessandra Mattei、Fabrizia Mealli、Shu Yang以及Anqi Zhao。在这本书中，我将分享从他们身上学到的知识。 


加州大学伯克利分校的许多学生对我早期版本的讲义提出了许多重要且富有建设性的意见。作为我“因果推断”课程的助教，艾米丽·弗拉纳根与司竹·陆认真阅读了我书籍的早期版本，并帮助我大幅改进了这本书。 


Joe Blitzstein教授仔细审阅了本书的早期版本，并给出了极为详尽的评论。对这些评论的采纳切实显著提升了本书的质量。Hongyuan Cao教授与Zhichao Jiang教授也基于本书的早期版本开设了“因果推断”课程，他们提出了许多极具价值的建议。 


我也要衷心感谢Min Young Woong、苏方舟、于超然以及袁洛华所提供的建议。


如果您发现任何错误，请随时给我发邮件。
